% -----------------------------------------------------------------------------
\subsubsection{Couleurs et lisibilité (F11)}
\label{sec:affichage}

L'application emprunte à l'application mobile Tado son principe~: la couleur d'une pièce dit
sa température. Un utilisateur de Tado ne doit pas être dépaysé.

\begin{exigence}{TPB-AFF-001}{Couleur selon la température}[
    justification = {Repérer les pièces froides et chaudes sans lire un chiffre.},
    priorite      = {Importante},
    verification  = {D},
    phase         = {Affichage},
    note          = {Quatre tranches~: moins de 18~°C (bleu), de 18 à 19,5~°C (cyan), de
                     19,5 à 21~°C (ambre), 21~°C et plus (orange).},
  ]
  Sur les montres couleur, l'application doit présenter chaque pièce de la liste dans une
  couleur qui dépend de sa température mesurée.
\end{exigence}

\begin{exigence}{TPB-AFF-002}{Palette de la montre}[
    justification = {Une couleur hors palette est approchée par la montre, et change.},
    priorite      = {Importante},
    verification  = {I},
    phase         = {Affichage},
  ]
  L'application ne doit employer que les couleurs de la palette de 64 couleurs des montres
  Pebble.
\end{exigence}

\begin{exigence}{TPB-AFF-003}{Jamais la couleur seule}[
    justification = {Les montres noir et blanc n'ont pas de couleurs~; une couleur seule
                     exclut aussi ceux qui distinguent mal les teintes.},
    priorite      = {Essentielle},
    verification  = {I},
    phase         = {Affichage},
  ]
  L'application ne doit porter aucune information par la couleur seule~: la température,
  l'état de chauffe, la sélection et l'intensité de l'historique doivent se lire aussi par
  un texte, une forme, un contour ou un motif.
\end{exigence}

\begin{exigence}{TPB-AFF-004}{Sélection visible}[
    justification = {Les boutons agissent sur la ligne choisie~; elle doit se voir de loin.},
    priorite      = {Essentielle},
    verification  = {D},
    phase         = {Affichage},
  ]
  L'application doit faire ressortir l'élément sélectionné de tous les autres.
\end{exigence}

\begin{exigence}{TPB-AFF-005}{Unités et formats}[
    justification = {Lire la température et l'heure comme chez soi.},
    priorite      = {Importante},
    verification  = {T},
    phase         = {Affichage},
    note          = {Les seuils de couleur (\refexig{TPB-AFF-001}) sont donnés en degrés
                     Celsius et convertis en degrés Fahrenheit.},
  ]
  L'application doit exprimer les températures dans l'unité du compte Tado, degré Celsius
  ou degré Fahrenheit, et les heures dans le format de la montre.
\end{exigence}

\begin{exigence}{TPB-AFF-006}{Langue}[
    justification = {L'utilisateur lit son chauffage dans sa langue.},
    priorite      = {Importante},
    verification  = {I},
    phase         = {Affichage},
    note          = {Le français et l'anglais sont pris en charge. Le séparateur décimal
                     suit la langue (virgule en français, point en anglais).},
  ]
  L'application doit présenter tous ses textes dans la langue de la montre, en français ou
  en anglais, et en anglais pour toute autre langue.
\end{exigence}

% -----------------------------------------------------------------------------
\subsubsection{Navigation (F12)}
\label{sec:navigation}

Une montre Pebble se manipule avec quatre boutons~: \emph{Haut}, \emph{Sélection},
\emph{Bas} et \emph{Retour}. Certaines montres ont aussi un écran tactile~: le toucher
est un complément facultatif des boutons, jamais un passage obligé. Il n'y a pas de clavier.
La figure~\ref{fig:navigation} montre l'enchaînement des écrans.

\begin{exigence}{TPB-NAV-001}{Ouverture sur la liste des pièces}[
    justification = {C'est la question la plus fréquente.},
    priorite      = {Essentielle},
    verification  = {D},
    phase         = {Navigation},
  ]
  À son ouverture, l'application doit afficher la liste des pièces.
\end{exigence}

\begin{exigence}{TPB-NAV-002}{Fonction des boutons}[
    justification = {Les montres Pebble se manipulent aux boutons, avec toujours le même
                     sens.},
    priorite      = {Essentielle},
    verification  = {D},
    phase         = {Navigation},
  ]
  Dans toutes les listes, l'application doit faire défiler la sélection avec les boutons
  \emph{Haut} et \emph{Bas}, ouvrir l'élément sélectionné avec \emph{Sélection} et revenir à
  l'écran précédent avec \emph{Retour}.
\end{exigence}

\begin{exigence}{TPB-NAV-003}{Menu général}[
    justification = {Les écrans qui concernent toute la maison ne sont pas des pièces.},
    priorite      = {Essentielle},
    verification  = {D},
    phase         = {Navigation},
  ]
  L'application doit ouvrir, par un appui long sur \emph{Sélection} depuis la liste des
  pièces, un menu qui donne accès à l'historique, à la présence, aux raccourcis et aux
  alertes.
\end{exigence}

\begin{exigence}{TPB-NAV-004}{Bien-être depuis le détail}[
    justification = {Même geste pour la même idée~: l'appui long mène à ce qui complète
                     l'écran.},
    priorite      = {Importante},
    verification  = {D},
    phase         = {Navigation},
  ]
  L'application doit ouvrir le bien-être ambiant de la pièce par un appui long sur
  \emph{Sélection} depuis le détail d'une pièce.
\end{exigence}

\begin{exigence}{TPB-NAV-005}{Quitter}[
    justification = {Les montres Pebble se quittent par \emph{Retour}.},
    priorite      = {Essentielle},
    verification  = {D},
    phase         = {Navigation},
  ]
  Depuis la liste des pièces, le bouton \emph{Retour} doit quitter l'application.
\end{exigence}

\begin{exigence}{TPB-NAV-006}{Tout faisable aux boutons}[
    justification = {Toutes les montres n'ont pas d'écran tactile, et le toucher n'est
                     pas toujours praticable (gants, mains mouillées).},
    priorite      = {Essentielle},
    verification  = {D},
    phase         = {Navigation},
  ]
  L'application doit permettre d'accomplir chacune de ses fonctions avec les seuls quatre
  boutons, que la montre ait ou non un écran tactile.
\end{exigence}

\begin{exigence}{TPB-NAV-007}{Toucher en complément}[
    justification = {Sur une montre tactile, toucher ce que l'on veut atteindre est plus
                     direct que d'appuyer plusieurs fois sur un bouton.},
    priorite      = {Souhaitable},
    verification  = {D},
    phase         = {Navigation},
    note          = {Fonction facultative, mise en \oe uvre écran par écran. Elle s'ajoute aux
                     boutons (\refexig{TPB-NAV-002}, \refexig{TPB-NAV-006}) et ne les remplace
                     jamais.},
  ]
  Sur les montres à écran tactile, l'application devrait permettre de manipuler ses écrans
  au toucher, en complément des boutons.
\end{exigence}

\begin{exigence}{TPB-NAV-008}{Liste des pièces au toucher}[
    justification = {La liste des pièces est l'écran le plus consulté~; la parcourir et
                     ouvrir une pièce d'un geste est le premier gain du toucher.},
    priorite      = {Souhaitable},
    verification  = {D},
    phase         = {Navigation},
  ]
  Sur les montres à écran tactile, l'application devrait permettre, dans la liste des
  pièces, de la faire défiler au toucher et d'ouvrir une pièce en la touchant, avec le même
  résultat que le bouton \emph{Sélection} sur cette pièce.
\end{exigence}

\begin{figure}[H]
  \centering
  \begin{tikzpicture}[
      font=\sffamily\footnotesize,
      ecr/.style={draw=exigBleu, fill=exigBleu!5, rounded corners=3pt, align=center,
                  text width=22mm, minimum height=9mm},
      men/.style={draw=exigGris, fill=exigGris!6, rounded corners=3pt, align=center,
                  text width=22mm, minimum height=9mm},
      fleche/.style={-Stealth, draw=exigGris!90!black},
      lbl/.style={font=\sffamily\scriptsize, text=exigBleu!80!black, align=center},
    ]
    \node[ecr] (pie) at (0, 0)    {\textbf{Pièces}};
    \node[ecr] (det) at (3.4, 0)  {\textbf{Détail}};
    \node[ecr] (pil) at (6.8, 0)  {\textbf{Pilotage}};
    \node[ecr] (con) at (10.2, 0) {\textbf{Forcer}\\température};
    \node[ecr] (dur) at (13.6, 0) {\textbf{Durée}};
    \node[ecr] (amb) at (3.4, -1.8) {\textbf{Bien-être}\\ambiant};
    \node[men] (men) at (0, -3.6)   {\textbf{Menu}};
    \node[ecr] (his) at (3.4, -3.6) {\textbf{Historique}};
    \node[ecr] (pre) at (6.8, -3.6) {\textbf{Présence}};
    \node[ecr] (rac) at (10.2, -3.6) {\textbf{Raccourcis}};
    \node[ecr] (ale) at (13.6, -3.6) {\textbf{Alertes}};
    \draw[fleche] (pie) -- node[lbl, above=1pt]{Sél.} (det);
    \draw[fleche] (det) -- node[lbl, above=1pt]{Sél.} (pil);
    \draw[fleche] (pil) -- node[lbl, above=1pt]{Sél.} (con);
    \draw[fleche] (con) -- node[lbl, above=1pt]{valider} (dur);
    \draw[fleche] (det) -- node[lbl, right=1pt]{Sél. long} (amb);
    \draw[fleche] (pie) -- node[lbl, left=1pt]{Sél. long} (men);
    \draw[fleche] (men) -- (his);
    \draw[fleche] (men.south) -- ++(0,-0.7) -| (pre.south);
    \draw[fleche] (men.south) -- ++(0,-0.7) -| (rac.south);
    \draw[fleche] (men.south) -- ++(0,-0.7) -| (ale.south);
  \end{tikzpicture}
  \caption{Enchaînement des écrans~; \emph{Retour} revient à l'écran précédent depuis chacun
    d'eux}
  \label{fig:navigation}
\end{figure}

% -----------------------------------------------------------------------------
\subsubsection{Modèles de montre (F13)}
\label{sec:modeles}

L'application est conçue pour la Pebble Time~2, et doit rester utilisable sur les autres
modèles Pebble. Les écrans diffèrent par leur taille, leur forme et leur nombre de couleurs~;
les fonctions, non.

\begin{table}[H]
  \centering\small
  \begin{tabularx}{\linewidth}{@{}>{\raggedright\arraybackslash}p{40mm}l>{\raggedright\arraybackslash}p{27mm}X@{}}
    \toprule
    \textbf{Modèle} & \textbf{Plateforme} & \textbf{Écran} & \textbf{Particularités} \\
    \midrule
    Pebble Time~2 & emery & 200\,×\,228, couleur & Barre d'actions de 30~px. Modèle de référence. \\
    Pebble Time, Time Steel & basalt & 144\,×\,168, couleur & Barre d'actions de 24~px. \\
    Pebble Time Round & chalk & 180\,×\,180, rond, couleur & Trois lignes, sélection au centre, pas de barre d'actions. Historique sur huit heures. \\
    Pebble Round~2 & gabbro & 260\,×\,260, rond, couleur & Comme chalk, avec les textes plus lisibles de emery. \\
    Pebble~2 & diorite & 144\,×\,168, noir et blanc & Sans couleur~: inversions, contours et motifs. \\
    Pebble~2 Duo & flint & 144\,×\,168, noir et blanc & Rendu identique à diorite. \\
    \bottomrule
  \end{tabularx}
  \caption{Modèles pris en charge}
  \label{tab:modeles}
\end{table}

\begin{exigence}{TPB-PLT-001}{Six plateformes}[
    justification = {Les montres Pebble existent en plusieurs plateformes d'affichage~; six sont retenues.},
    priorite      = {Essentielle},
    verification  = {D},
    phase         = {Modèles},
    note          = {Les Pebble et Pebble Steel (plateforme aplite) ne sont pas prises en charge~: leur
                     mémoire de travail de 24~Ko ne laisse que 3~Ko de libre une fois le code chargé
                     (décision du 29 octobre 2026).},
  ]
  L'application doit fonctionner sur les montres des six plateformes du
  tableau~\ref{tab:modeles}.
\end{exigence}

\begin{exigence}{TPB-PLT-002}{Mêmes fonctions partout}[
    justification = {Choisir sa montre ne doit pas priver de fonctions.},
    priorite      = {Essentielle},
    verification  = {D},
    phase         = {Modèles},
  ]
  L'application doit offrir sur chaque plateforme les mêmes écrans et les mêmes fonctions,
  l'adaptation ne portant que sur leur présentation.
\end{exigence}

\begin{exigence}{TPB-PLT-003}{Noir et blanc}[
    justification = {Les Pebble et Pebble~2 n'ont pas de couleur.},
    priorite      = {Essentielle},
    verification  = {D},
    phase         = {Modèles},
  ]
  Sur les montres noir et blanc, l'application doit présenter une pièce qui chauffe en
  vidéo inversée, la sélection par un contour épais et les niveaux de l'historique par des
  motifs distincts.
\end{exigence}

\begin{exigence}{TPB-PLT-004}{Écrans ronds}[
    justification = {Sur un écran rond, les lignes du haut et du bas sont rognées.},
    priorite      = {Essentielle},
    verification  = {D},
    phase         = {Modèles},
  ]
  Sur les montres à écran rond, l'application doit présenter les listes sur trois lignes
  visibles, l'élément sélectionné occupant la ligne du centre.
\end{exigence}

\begin{exigence}{TPB-PLT-005}{Barre d'actions}[
    justification = {Les écrans rectangulaires ont la place d'afficher les gestes
                     possibles~; les écrans ronds non.},
    priorite      = {Souhaitable},
    verification  = {D},
    phase         = {Modèles},
  ]
  Sur les montres à écran rectangulaire, l'application devrait afficher à côté des boutons
  les actions qu'ils déclenchent sur l'écran en cours.
\end{exigence}

\begin{exigence}{TPB-PLT-006}{Textes abrégés}[
    justification = {\enquote{Salle de bain} ne tient pas sur une ligne de 144 pixels.},
    priorite      = {Importante},
    verification  = {D},
    phase         = {Modèles},
  ]
  Lorsque la place manque, l'application doit abréger les noms de pièce et les
  indications, sans que deux pièces deviennent indiscernables.
\end{exigence}

\begin{exigence}{TPB-PLT-007}{Durée de l'historique selon l'écran}[
    justification = {Douze colonnes ne tiennent pas dans la largeur utile d'un écran rond de
                     180 pixels.},
    priorite      = {Importante},
    verification  = {D},
    phase         = {Modèles},
  ]
  Sur la Pebble Time Round, l'application doit présenter l'historique des huit dernières
  heures, au lieu des douze.
\end{exigence}

